
% ─── Attori ───────────────────────────────────────────────────
\section*{Attori}

\begin{longtable}{L{3.5cm} L{8.5cm} L{3cm}}
\toprule
\textbf{Termine} & \textbf{Descrizione} & \textbf{Sinonimi} \\
\midrule
\endhead
\bottomrule
\endfoot

Chef
& Stabilisce il \textbf{menù} per gli \textbf{eventi} e ne supervisiona
  la preparazione.
& — \\
\addlinespace

Cliente
& Colui che commissiona l'organizzazione di un \textbf{evento}.
& — \\
\addlinespace

Cuoco
& Prepara il \textbf{cibo}.
& — \\
\addlinespace

Organizzatore
& La persona che gestisce il \textbf{personale} e gli \textbf{eventi}.
& — \\
\addlinespace

Personale
& La parola \emph{personale} è usata in modo ambiguo nel testo per intendere
  talvolta tutti i dipendenti (\textbf{chef}, \textbf{cuochi},
  \textbf{personale di servizio}), talvolta soltanto i dipendenti soggetti
  a \textbf{turni} (\textbf{cuochi} e \textbf{personale di servizio}).
  In presenza di questo termine consultare il cliente per disambiguare.
& — \\
\addlinespace

Personale di Servizio
& Le persone (maître e camerieri) che si occupano del \textbf{servizio}
  durante l'\textbf{evento} stesso.
& Staff di supporto \\
\addlinespace

Utente
& Generico utilizzatore dell'applicazione (ha necessariamente un ruolo fra
  \textbf{organizzatore}, \textbf{chef}, \textbf{cuoco},
  \textbf{personale di servizio}).
& — \\

\end{longtable}

% ─── Dominio ──────────────────────────────────────────────────
\section*{Dominio dell'Applicazione}

\begin{longtable}{L{3.5cm} L{8.5cm} L{3cm}}
\toprule
\textbf{Termine} & \textbf{Descrizione} & \textbf{Sinonimi} \\
\midrule
\endhead
\bottomrule
\endfoot

Annullamento di un evento
& Un \textbf{evento} può essere annullato in qualsiasi fase della sua
  predisposizione. Se l'annullamento avviene a ridosso della data
  dell'\textbf{evento} (tipicamente a meno di una settimana), al
  \textbf{cliente} viene addebitata una \textbf{penale}.
& — \\
\addlinespace

Autore (di una ricetta)
& Chi ha ideato la \textbf{ricetta}.
& — \\
\addlinespace

Capofila (di un evento ricorrente)
& L'\textbf{evento ricorrente} da cui le singole \textbf{istanze}
  ereditano le impostazioni di default (luogo, durata, numero di
  servizi richiesti, numero di partecipanti, note e cliente). Le modifiche al
  capofila possono essere propagate a tutte le istanze ancora
  modificabili; ogni \textbf{istanza} può tuttavia essere modificata
  individualmente senza alterare le altre.
& — \\
\addlinespace

Cibo
& Piatti preparati (seguendo una \textbf{ricetta} del \textbf{ricettario})
  per essere consumati durante un \textbf{servizio}.
& Pietanze \\
\addlinespace

Compito (di cucina)
& Attività assegnata dallo \textbf{chef} ai \textbf{cuochi} durante la fase
  di \textbf{lavoro preparatorio}, che consiste nello svolgimento di uno o
  più \textbf{passi} di una \textbf{ricetta} o \textbf{preparazione} volti
  alla realizzazione di un \textbf{piatto} o di un \textbf{preparato
  intermedio} in vista di un \textbf{servizio} che deve essere effettuato.
& — \\
\addlinespace

Conclusione (di un evento ricorrente)
& La data fino a cui l'\textbf{evento ricorrente} si ripete, oppure un
  numero massimo di ripetizioni che avranno luogo.
& — \\
\addlinespace

Disponibilità
& Indicazione fornita da un \textbf{cuoco} o dal \textbf{personale di
  servizio} riguardo la propria presenza per un \textbf{turno}. La
  disponibilità può essere revocata finché la persona non viene
  effettivamente convocata/assegnata al turno da \textbf{chef} od
  \textbf{organizzatore}, momento in cui diventa un vincolo non più
  ritirabile.
& — \\
\addlinespace

Documentazione
& Informazioni e allegati inseriti dall'\textbf{organizzatore} al momento
  del passaggio dell'\textbf{evento} allo stato \emph{terminato}.
& Note di chiusura \\
\addlinespace

Dose (di un ingrediente)
& Indica la quantità di ogni \textbf{ingrediente} (di base o preparati)
  necessario alla preparazione di un \textbf{piatto}.
& — \\
\addlinespace

Eliminazione di un evento
& Rimozione definitiva della \textbf{scheda di un evento} e di tutti i suoi
  dati. Un \textbf{evento} può essere eliminato solo quando si trova in
  stato di \emph{bozza}, ossia prima che i \textbf{menù} vengano
  approvati e i lavori abbiano inizio. Una volta che l'\textbf{evento}
  è \emph{in corso} non è più eliminabile, ma può essere
  \textbf{annullato} (si veda \emph{annullamento di un evento}).
& — \\
\addlinespace

Evento
& Contesto in cui viene fornito il \textbf{servizio} di catering. Prevede due
  momenti diversi: il \textbf{lavoro preparatorio} e il \textbf{servizio}.
  Se ne fa carico un \textbf{organizzatore}, mentre la cucina è affidata a
  uno \textbf{chef}.
  Un evento può essere \emph{semplice}, e prevedere un singolo \textbf{servizio},
  o \emph{complesso}, prevedendo più \textbf{servizi} (pranzo e cena, colazione
  e pranzo, coffee-break mattino e pomeriggio, ecc.) in una o più giornate.
  Un evento attraversa i seguenti stati:
  \begin{itemize}[nosep,leftmargin=*]
    \item \emph{bozza}: l'\textbf{organizzatore} sta definendo la
          \textbf{scheda di un evento}; l'evento è eliminabile;
    \item \emph{in corso}: i \textbf{menù} sono stati approvati e i lavori
          hanno avuto inizio; l'evento non è più eliminabile, ma può essere
          annullato;
    \item \emph{terminato}: l'\textbf{organizzatore} ha chiuso l'evento al
          termine del \textbf{servizio};
    \item \emph{annullato}: l'evento è stato interrotto prima della sua
          conclusione naturale (si veda \textbf{annullamento di un evento}).
  \end{itemize}
& — \\
\addlinespace

Evento ricorrente
& \textbf{Evento} che si ripete con una data regolarità. In particolare è
  rilevante:
  \begin{itemize}[nosep,leftmargin=*]
    \item ogni quanto l'evento si ripete (\textbf{frequenza});
    \item la data fino a cui si ripete, o il numero di occorrenze
          (\textbf{conclusione}).
  \end{itemize}
& — \\
\addlinespace

Frequenza di un evento ricorrente
& Intervallo di tempo dopo cui l'\textbf{evento ricorrente} si ripete.
& — \\
\addlinespace

Gestione del servizio
& La gestione dei \textbf{servizi} che compongono un \textbf{evento} viene
  affidata a un \textbf{organizzatore}. Essa prevede:
  \begin{itemize}[nosep,leftmargin=*]
    \item scegliere il \textbf{personale di servizio} per ogni
          \textbf{turno di servizio} associato all'\textbf{evento},
          indicando il \textbf{ruolo} che avrà in quella particolare
          situazione;
    \item opzionalmente, l'assegnamento di un \textbf{cuoco} a quello stesso
          \textbf{servizio} se il \textbf{menù} prevede \textbf{ricette}
          che lo rendono necessario.
  \end{itemize}
& — \\
\addlinespace

Ingrediente
& Materia prima per preparare un \textbf{piatto}. Può essere \emph{di base}
  (scelto da un elenco di partenza corrispondente a ciò che si trova in
  dispensa) oppure un \textbf{preparato intermedio} realizzato in cucina
  dagli stessi \textbf{cuochi}.
& — \\
\addlinespace

Insieme delle istruzioni (di una ricetta)
& Spiegazione di come trasformare gli \textbf{ingredienti} di partenza in un
  prodotto finito. Sono divise in due sezioni: la parte realizzabile in
  anticipo e quella da realizzare all'ultimo sul posto dell'evento. Ogni
  singola istruzione è rappresentata da un \textbf{passo}.
& — \\
\addlinespace

Istanza (di evento)
& Singola edizione di un \textbf{evento ricorrente}. Di default
  eredita le impostazioni del \textbf{capofila}, ma può essere
  modificata individualmente senza influenzare le altre istanze
  della ricorrenza.
& Occorrenza \\
\addlinespace

Lavoro di servizio
& Momento di lavoro che si svolge nel luogo dell'\textbf{evento}, effettuato
  dal \textbf{personale di servizio} e opzionalmente da \textbf{cuochi} di
  supporto. È organizzato in \textbf{turni di servizio}.
& — \\
\addlinespace

Lavoro preparatorio
& Momento di lavoro che si svolge in \textbf{sede}, effettuato dai
  \textbf{cuochi} sotto la supervisione di uno \textbf{chef}. È organizzato
  in \textbf{turni preparatori}.
& Preparazione in sede \\
\addlinespace

Menù
& Descrive quali piatti vengono proposti in un dato \textbf{servizio}
  (nell'ambito di un \textbf{evento}). Si compone di diverse \textbf{voci},
  opzionalmente organizzate in \textbf{sezioni}. Lo \textbf{chef} costruisce
  i suoi menù a partire dalle \textbf{ricette} del \textbf{ricettario}.
  Un menù è caratterizzato da informazioni aggiuntive:
  \begin{itemize}[nosep,leftmargin=*]
    \item se è consigliata la presenza di un \textbf{cuoco} durante il
          \textbf{servizio};
    \item se prevede solo piatti freddi o anche piatti caldi;
    \item se richiede la disponibilità di una cucina nella \textbf{sede}
          dell'\textbf{evento};
    \item se è adeguato per un buffet;
    \item se può essere fruito senza posate (finger food).
  \end{itemize}
  Un menù può essere modificato fintanto che non è utilizzato in alcun
  \textbf{evento}.
& — \\
\addlinespace

Modifica al menù (di un evento)
& Variazione temporanea proposta dall'\textbf{organizzatore} prima
  dell'approvazione del \textbf{menù} per un \textbf{evento} (aggiungendo o
  togliendo piatti). Tali modifiche non alterano il \textbf{menù} originale
  presente nel \textbf{ricettario}, restano visibili come aggiunte o
  eliminazioni limitate all'\textbf{evento} in questione e devono essere
  confermate dallo \textbf{chef}.
& Proposta di modifica al menù \\
\addlinespace

Passo
& Singolo elemento che compone le istruzioni di una \textbf{ricetta}. I passi
  si dividono in due sezioni: quella dei passi realizzabili in anticipo
  (durante il \textbf{lavoro preparatorio}) e quella dei passi da effettuare
  sul posto (durante il \textbf{servizio}).
& Istruzione \\
\addlinespace

Penale
& Sanzione economica addebitata al \textbf{cliente} in caso di
  \textbf{annullamento di un evento} tardivo (tipicamente a meno di una
  settimana dall'evento) oppure di modifica sostanziale del numero di
  partecipanti dopo l'avvio dei lavori. L'\textbf{organizzatore} può, a
  propria discrezione, concedere deroghe (esenzioni) dalla penale a
  clienti affezionati.
& — \\
\addlinespace

Piatto
& Risultato dell'esecuzione di una \textbf{ricetta}, che può essere servito al
  \textbf{cliente} in quanto nel suo stato finale.
& — \\
\addlinespace

Porzione
& Quantità indicativa di un \textbf{piatto} finito per una persona.
& — \\
\addlinespace

Preparato intermedio
& Composto realizzato tramite le indicazioni di una \textbf{preparazione}, che
  diventa \textbf{ingrediente} di un'altra \textbf{preparazione} o di una
  \textbf{ricetta}.
& Preparato di base \\
\addlinespace

Preparazione (nel ricettario)
& Del tutto analoga a una \textbf{ricetta}, ma descrive come realizzare un
  \textbf{preparato intermedio} da utilizzare in un'altra \textbf{preparazione}
  o \textbf{ricetta}. Attraversa gli stessi stati di una \textbf{ricetta}:
  \begin{itemize}[nosep,leftmargin=*]
    \item \emph{bozza}: visibile solo al creatore; modificabile ed
          eliminabile;
    \item \emph{pubblicata}: visibile a tutti; non più modificabile se in
          uso, ritirabile solo se non utilizzata.
  \end{itemize}
& — \\
\addlinespace

Procedura (di preparazione)
& Sequenza di \textbf{passi} indicata da una \textbf{ricetta} (o
  \textbf{preparazione}) che un \textbf{cuoco} deve eseguire per realizzare
  il \textbf{piatto} o il \textbf{preparato intermedio}.
& — \\
\addlinespace

Proprietario (di una ricetta)
& Chi ha inserito la \textbf{ricetta} nel \textbf{ricettario}.
& — \\
\addlinespace

Raggruppamento di turni
& Insieme di uno o più \textbf{turni preparatori} per cui il
  \textbf{personale} fornisce un'unica \textbf{disponibilità}, valida per
  tutti i turni del raggruppamento o per nessuno; le regole di modifica
  della disponibilità si applicano all'intero raggruppamento. Può essere
  \emph{singolo} (raggruppa due o più turni specifici una tantum) o \emph{ricorrente}
(raggruppa turni che si ripetono con una data periodicità nel tempo).
& — \\
\addlinespace

Ricetta
& Descrive come preparare un \textbf{piatto} da servire a tavola a partire
  dagli \textbf{ingredienti} (ciascuno con una propria \textbf{dose}) e dalle
  istruzioni per realizzarla. È caratterizzata da un nome, da un
  \textbf{proprietario}, opzionalmente da un \textbf{autore}, e può essere
  accompagnata da una descrizione breve. Ha inoltre un tempo di preparazione
  e può essere accompagnata da tag (vegetariana, senza latticini, senza uova,
  senza glutine, vegana).
  Produce un \textbf{piatto} finito (a differenza di una \textbf{preparazione},
  che produce un \textbf{preparato intermedio}).
  Attraversa i seguenti stati:
  \begin{itemize}[nosep,leftmargin=*]
    \item \emph{bozza}: visibile solo al creatore; modificabile ed
          eliminabile;
    \item \emph{pubblicata}: visibile a tutti; non più modificabile se in
          uso, ritirabile solo se non utilizzata.
  \end{itemize}
& — \\
\addlinespace

Ricettario
& Un elenco/raccolta di \textbf{ricette} e \textbf{preparazioni}.
& — \\
\addlinespace

Richiesta (di organizzazione)
& Un \textbf{cliente} commissiona l'organizzazione di un \textbf{evento}
  all'azienda di catering.
& — \\
\addlinespace

Ruolo
& Specifica mansione assegnata a un membro del \textbf{personale di servizio}
  in un dato \textbf{servizio}.
& — \\
\addlinespace

Scheda di un evento
& Creata dall'\textbf{organizzatore}, è il documento in cui sono descritte le
  caratteristiche dell'\textbf{evento}. Riporta informazioni quali:
  \begin{itemize}[nosep,leftmargin=*]
    \item luogo;
    \item date;
    \item durata;
    \item numero di persone;
    \item numero di \textbf{servizi};
    \item \textbf{cliente}/committente;
    \item note particolari (opzionale).
  \end{itemize}
& — \\
\addlinespace

Sede
& Ubicazione della società di catering e delle relative cucine.
& — \\
\addlinespace

Sede di un evento
& Luogo di svolgimento degli \textbf{eventi}.
& — \\
\addlinespace

Servizio
& Attività che consiste nel servire ai partecipanti i piatti e le bevande
  preparati. Include l'allestimento preliminare e il rigoverno al termine.
  Può andare dal semplice buffet all'allestimento di una sala ristorante.
  Uno stesso \textbf{evento} può prevedere più servizi, anche in
  \textbf{sedi} diverse. Ciascun servizio avrà una precisa fascia oraria,
  e un proprio \textbf{menù} e un proprio \textbf{personale di servizio}.
  Un singolo servizio appartenente a un \textbf{evento} può essere
  annullato indipendentemente dagli altri servizi, assumendo lo stato di
  \emph{annullato} per permettere di liberare il personale assegnato.
& — \\
\addlinespace

Sezione di un menù
& Partizionamento (opzionale) di un \textbf{menù}. Ad esempio: primi,
  secondi, dolci sono sezioni.
& — \\
\addlinespace

Turno
& Periodo di attività di una persona caratterizzato da una data, un luogo di
  svolgimento e da una fascia oraria. I turni sono di due tipi:
  \textbf{turni preparatori} e \textbf{turni di servizio}. Un turno può
  essere indicato come ``completo'' se la cucina e lo staff per quel turno
  sono completamente impegnati.
& — \\
\addlinespace

Turno preparatorio
& \textbf{Turno} in cui i \textbf{cuochi} lavorano a preparare le pietanze.
  I turni preparatori possono essere raggruppati in un
  \textbf{raggruppamento di turni}: il \textbf{personale} dà la propria
  \textbf{disponibilità} per tutti i turni del raggruppamento o per nessuno.
& Turno di cucina \\
\addlinespace

Turno di servizio
& \textbf{Turno} in cui lavorano il \textbf{personale di servizio} e i
  \textbf{cuochi} per realizzare il \textbf{servizio} stesso. Può
  specificare un tempo aggiuntivo rispetto agli orari del \textbf{servizio}
  effettivo per la preparazione (prima) e per rigovernare (dopo).
& — \\
\addlinespace

Voce di un menù
& Singolo elemento di un \textbf{menù}. Può appartenere a una \textbf{sezione}
  se il \textbf{menù} è diviso in sezioni. Ogni voce fa riferimento a una
  \textbf{ricetta} nel \textbf{ricettario}, ma il testo della voce può essere
  diverso dal nome della \textbf{ricetta}.
& — \\

\end{longtable}
